%% ---------------------------------------------------------------------------
%% Revised Section 5 (results).
%% ---------------------------------------------------------------------------

\section{Experiments and Results} \label{discussion}

This section applies the protocol of Section~\ref{methodology} to two benchmarks: Section~\ref{configuration} gives the setup and the teacher--student baseline, Section~\ref{sec:dim1} the four evaluation dimensions, Section~\ref{sec:logistic} error detection against baselines, and Section~\ref{sec:calibration} what the surrogate inherits from its teacher. Unless stated otherwise, ``error'' means an error with respect to the gold label (teacher disagreement: Table~S5).

\subsection{Experimental Setup and Teacher--Student Baseline} \label{configuration} \label{sec:llm_teacher}

We use AG News~\cite{10.5555/2969239.2969312}, a four-class topic classification task with 120,000 training and 7,600 test instances, and SST-2~\cite{socher-etal-2013-recursive}, binary sentiment classification with 67,349 training instances and the 872 labelled instances of its validation split as test set. The language model is BERT, fine-tuned once per task; the surrogate is a GCN with a multilayer perceptron classifier head, trained on the fine-tuned model's embeddings and predictions (Section~\ref{gnn}) on the full training split of each dataset. The checkpoint with the lowest loss on the test split is kept, with early stopping after five epochs without improvement, and the error-detection analyses are cross-validated on that same split. Every explainer runs under the common budget of Section~\ref{sec:crossfair} (2000 forward passes and a 20\% element budget per instance). Four graph representations per dataset plus the token representation give 18 explainer--representation configurations.



Table~\ref{tab:performance_comparison} reports precision, recall and F\(_1\) of the fine-tuned BERT teacher and of each GCN surrogate, against the teacher's predictions and against the gold labels. With node features taken from the fine-tuned teacher, agreement with the teacher is 97--98\% in every configuration; against the gold labels the surrogates are within 0.3 points of BERT on AG News and 0.3--1.3 points below it on SST-2. Most of the surrogates' gold-label errors are errors of the teacher itself (Section~\ref{sec:calibration}); the gap comes from the 1.7--3.1\% of test instances on which surrogate and teacher disagree, where the teacher is right more often than the surrogate (Section~\ref{sec:logistic}), and the larger gap on SST-2 lies within the range of the bootstrap resampling of its test instances (Section~\ref{logistic_regression}): the 95\% interval for the gap covers the AG News values for every graph type.

%% Table 2 (tab:performance_comparison) unchanged.

\subsection{The Four Evaluation Dimensions} \label{sec:dim1} \label{sec:dim2} \label{sec:dim3} \label{sec:dim4}

\subsubsection{Dimension 1: AUC discriminative capacity}

Deletion AUC integrates the confidence in the predicted class as the elements of the explanation are removed one at a time in decreasing order of importance, the most important first, until the whole explanation is removed; insertion AUC integrates it as the same elements are added to an empty input in the same order. Both curves therefore span the 20\% budget rather than the whole input.

Figure~\ref{fig:thresholds} gives the detection rates by threshold and Table~S11 the AUROC of each AUC used alone as an error score. The deletion AUC is the more discriminative of the two in every configuration, and for every explainer it stays close to the maximum softmax probability (MSP) of the explained model, that is, its confidence in the predicted class (0.81--0.88 against 0.84--0.89). The deletion curve starts at the MSP and every later point is the confidence that survives the removal; when removing what the explanation highlights moves the model little within the budget (Table~\ref{tab:signatures}), the curve stays flat and its integral is essentially the MSP. The per-instance Spearman correlation between the two is 0.68--1.00, above 0.80 in 16 of the 18 configurations: the deletion AUC ranks the instances as the confidence does, which is why it detects errors almost as well as the confidence and says little about the explanation. The insertion curve starts from an empty input, where the confidence is unrelated to the MSP, and rises as the ranked elements are added, so its integral depends on the order in which the explainer ranks them. With the surrogate fixed, on the same graphs, SubgraphX and GraphSVX give the same deletion AUC (means within 0.04, per-instance Spearman 0.75--0.99 between the two explainers) but insertion AUCs that differ by up to 0.35 in the mean and rank the instances differently (Spearman 0.45--0.92). The first is a property of the model, the second of the explanation. Not carrying the confidence, the insertion AUC is the weaker error signal for every explainer (0.56--0.82).

\begin{figure}[H]
    \centering
    \includegraphics[width=\linewidth]{figures/Figure_3_thresholds.png}
    \caption{Error and correctness detection rates when a prediction is flagged as an error because its insertion AUC (top) or deletion AUC (bottom) falls below a threshold, pooled over the four graph types of each explainer (gold-label errors, fine-tuned node features). The deletion curves of the three explainers coincide; the insertion curves differ by explainer.}
    \label{fig:thresholds}
\end{figure}

\paragraph{\textbf{Key Finding}}
Under the 20\% budget, the deletion trajectory measures the model and the insertion trajectory measures the explanation: deletion is the model's confidence in another guise for every explainer and representation, insertion is what distinguishes explainers and the weaker error signal. No trajectory used alone improves on the confidence of the explained model.

\subsubsection{Dimension 2: feature ranking stability and importance concentration}

Dimension~2 reads the two trajectories of Dimension~1 at fixed points instead of integrating them. The mask-out drop at \(k\), \(c(x)-c(x_{\setminus S_k})\), is the confidence lost when the \(k\) top-ranked elements are removed; the sufficiency drop, \(c(x)-c(x_{S_k})\), is the confidence still missing when only those \(k\) elements are kept, for \(k\in\{1,3,5,10\}\) (the last element of the explanation when it has fewer than \(k\)); the concentration at \(k\) is the share of the drop over the whole explanation that the first \(k\) elements account for. Two things change with respect to Dimension~1: the reading is at a fixed number of elements, a common granularity for all graph sizes, and the original confidence is subtracted, so the metric does not carry the MSP as the deletion AUC does. Figure~S1 plots the drops by explainer and correctness. For the graph explainers the mask-out drop at \(k=5\) separates correct from incorrect predictions (AUROC 0.74--0.84 on AG News, 0.58--0.81 on SST-2; Table~S11); for TokenSHAP it carries almost no information (0.51--0.52). The difference is one of direction. For every surrogate, removing its top-ranked nodes lowers the confidence, in the same direction for almost every instance, and by more on its errors than on its correct predictions. For the language model each removed word can push the confidence either way: on BERT's errors, removing the five top-ranked words raises it in 20--26\% of the cases (3--5\% of the correct predictions), so the signed drop cancels in a rank-based measure while its magnitude separates errors as well as the graph metrics do (AUROC 0.85 and 0.77).

\paragraph{\textbf{Key Finding}}
Concentration-based metrics carry information about correctness for both families, in the magnitude of the drop for the language model and in its magnitude and sign for the surrogates, but that information is the same that the model's confidence carries (Section~\ref{sec:logistic}).

\subsubsection{Dimension 3: consistency across outcomes}

Dimension~3 reads the same two perturbations on the decision margin, the difference between the probability of the predicted class and the highest probability among the other classes, on the same input: \(m(x_S)\) is the margin when only the explanation is kept and \(m(x_{\setminus S})\) the margin once it is removed. A margin that stays positive means the decision survives the perturbation; one that changes sign means the prediction flips.

Removing the explanation rarely changes the decision, whatever the explainer and the topology, and the margin that remains keeps the order of the original ones. Since errors are on average the least confident predictions (the confidence baseline of Figure~\ref{fig:error_detection}), \(m(x_{\setminus S})\) separates them from correct predictions as well as the original confidence does, and it is, with the deletion AUC, the most discriminative explanation metric for every explainer.

The flips of Figure~\ref{fig:flips} concern at most 5\% of the correct predictions of the surrogates and 3--11\% of those of the language model, but up to 30\% and 20--32\% of their errors. The margin of the explanation alone, \(m(x_S)\), is a weaker error signal and separates the explainers rather than the topologies: it reproduces the decision nearly always for GraphSVX and less often for SubgraphX and TokenSHAP, for errors as well as for correct predictions, so it describes the explanation rather than signalling an error (per-configuration values in Tables~S10 and~S12).

\paragraph{\textbf{Key Finding}}
An error is far more likely than a correct prediction to flip to another class when its explanation is removed, and the tendency is strongest for the language model, which depends more on its top-ranked tokens than the surrogates on their top-ranked nodes.

\begin{figure}[H]
    \centering
    \includegraphics[width=\linewidth]{figures/Figure_4_flips.png}
    \caption{Dimension 3 read as decisions: percentage of predictions whose decision flips when the explanation is removed (\(m(x_{\setminus S})<0\)), for correct predictions (light) and errors (dark), by explainer and graph type (gold-label errors; values in Table~S10).}
    \label{fig:flips}
\end{figure}

\subsubsection{Dimension 4: behavioural faithfulness}

Dimension~4 measures what fraction of the original confidence the explanation accounts for: sufficiency \(F^{+}=c(x_S)/c(x)\) is the fraction that the explanation reproduces on its own, and necessity \(F^{-}=1-c(x_{\setminus S})/c(x)\) the fraction that is lost when the explanation is removed. Table~\ref{tab:signatures} lists both per explainer and topology.

\begin{table}[H]
\centering
\scriptsize
\setlength{\tabcolsep}{4pt}
\caption{Explanation signatures by explainer and topology (Dimension~4): mean sufficiency \(F^{+}=c(x_S)/c(x)\) and mean necessity \(F^{-}=1-c(x_{\setminus S})/c(x)\). Every configuration is sufficient (\(F^{+}\ge 0.5\)) and not necessary (\(F^{-}\le 0.12\)). The realised fraction of input elements in the explanation is 0.16--0.35 in every configuration.}
\label{tab:signatures}
\begin{tabular}{@{}llcccc@{}}
\toprule
& & \multicolumn{2}{c}{\textbf{AG News}} & \multicolumn{2}{c}{\textbf{SST-2}} \\
\cmidrule(lr){3-4} \cmidrule(lr){5-6}
\textbf{Explainer} & \textbf{Graph} & \(F^{+}\) & \(F^{-}\) & \(F^{+}\) & \(F^{-}\) \\
\midrule
\multirow{4}{*}{SubgraphX} & Constituency & 0.50 & 0.01 & 0.71 & 0.01 \\
 & Syntactic & 0.83 & 0.04 & 0.96 & 0.09 \\
 & Window & 0.57 & 0.01 & 0.81 & 0.02 \\
 & Skip-gram & 0.51 & 0.01 & 0.66 & 0.01 \\
\midrule
\multirow{4}{*}{GraphSVX} & Constituency & 0.92 & 0.04 & 0.98 & 0.03 \\
 & Syntactic & 0.98 & 0.12 & 0.99 & 0.11 \\
 & Window & 0.72 & 0.01 & 0.91 & 0.02 \\
 & Skip-gram & 0.78 & 0.01 & 0.93 & 0.02 \\
\midrule
TokenSHAP & Tokens & 0.76 & 0.02 & 0.67 & 0.11 \\

\end{tabular}
\end{table}




All 18 configurations fall in the same quadrant, sufficient (high \(F^{+}\)) and not necessary (low \(F^{-}\); Table~\ref{tab:signatures}): the elements the explainer selects reproduce on their own 50--99\% of the model's confidence, and the elements left outside the budget sustain 88--99\% of it, so the prediction rests on more evidence than the explanation contains. Within that quadrant the degrees differ: GraphSVX explanations are more sufficient than SubgraphX explanations on every one of the eight graph types, SubgraphX is least sufficient on the constituency graphs, where the connectivity constraint fills its coalition with phrase nodes (66--75\% of the coalition) that share the [CLS] feature and carry little class evidence, and TokenSHAP lies within the range of the graph explainers on both datasets.

\paragraph{\textbf{Key Finding}}
Within the budget every explanation is sufficient (high \(F^{+}\)) and almost none is necessary (low \(F^{-}\)), whatever the explainer and the representation: the input left outside the budget carries enough evidence to sustain the prediction, so removing the explanation costs a fraction of the confidence rather than the decision. Degrees of sufficiency differ between explainers, but they do not separate graph from token explanations.

\subsection{Error Detection Against Baselines} \label{sec:logistic}

Taken together, the four dimensions carry information about correctness. The question of this section is where that information comes from: whether it is a property of the explanation itself, or whether the explanation inherits it from the confidence of the model it explains.

Figure~\ref{fig:error_detection} answers it with three error detectors per configuration, each a class-balanced logistic regression (Section~\ref{logistic_regression}) trained on a different input: the four dimensions of explanation features, the model's own confidence (three label-free statistics of its output distribution), or both. The AUROC of the confidence detector is how well the model detects its own errors without any explanation. If the explanation carries information of its own, adding it to the confidence raises the AUROC; if it only inherits the confidence, the AUROC stays the same. The bootstrap interval of that difference says whether the change is real. The confidence of the language model is drawn as a reference. Accuracy is left to Table~S12: the trivial detector that declares every prediction correct is right on 90--94\% of the test instances and detects no error, and a class-balanced detector never exceeds it because it flags errors at the cost of some correct predictions.

\begin{figure}[H]
    \centering
    \includegraphics[width=0.95\linewidth]{figures/Figure_5_error_detection.png}
    \caption{Error detection with respect to the gold label (10-fold stratified cross-validation, class-balanced logistic regression): AUROC of the detector trained on the explanation features (squares), on the model's confidence (circles) and on both (diamonds, with the 95\% paired-bootstrap interval of the gain over the confidence), for every configuration; dashed line, the detector on BERT's confidence. Values in Table~S12.}
    \label{fig:error_detection}
\end{figure}

Used alone, the explanation features detect errors in every configuration, but modestly and to the same degree everywhere: AUROC 0.79--0.86, with at most 0.05 between explainers and topologies (balanced accuracy 0.71--0.79 and MCC 0.30--0.38, Table~S12). The language model's own explainer is the best detector on AG News and within 0.02 AUROC of the best graph configuration on SST-2, so no representation has an advantage at this stage.

On AG News the surrogate's confidence is clearly a weaker error signal than BERT's, by about twice the width of the bootstrap intervals, because the surrogate is over-confident in the teacher's errors (Section~\ref{sec:calibration}); adding the explanation recovers part of that loss, significantly in seven of the eight configurations, and never all of it. On SST-2 the surrogate's confidence is as good an error signal as BERT's, every difference between models lies within the intervals, and adding the explanation gains nothing. Added to BERT's own confidence, the explanation changes nothing on either dataset (Table~S12). The explanation therefore helps exactly where distillation has damaged the confidence, and only up to what was damaged; it adds nothing to a confidence that is intact.

No surrogate detector reaches the language model's confidence alone (dashed line), with a linear or a gradient-boosting detector (Table~S6); each detector detects its own model's errors, and the next paragraph compares on the same errors.

\paragraph{Combining the surrogate with the language model}
A natural objection is that the surrogate's signals should be combined with, not compared to, the language model's. We therefore add to the detector of \emph{BERT's} gold errors on BERT's own confidence the surrogate's confidence, its disagreement with BERT, its confidence and explanation features, and committees of the four topologies (Section~\ref{logistic_regression}). None of the 228 combinations improves on BERT's confidence in AUROC, and the committees that pool confidence and explanation features are significantly worse (\(\Delta\)AUROC \(-0.01\) to \(-0.02\) on AG News, \(-0.06\) to \(-0.10\) on SST-2). Five combinations give an AUPRC gain whose interval lies at or just above zero (lower bounds 0.000--0.003): the surrogate's disagreement on the proximity graphs of AG News and, among BERT's high-confidence predictions on SST-2, the signals of one SubgraphX surrogate (Table~S7).

The surrogate is not a useful second opinion either: on the 1.7--3.1\% of test instances where the two disagree, BERT is right more often (by a factor of 1.1--2.4), and any rule that defers to the surrogate lowers accuracy (Table~S7).

\paragraph{Dimension weights}
The coefficients of the combined detector tell which explanation features still matter once the confidence is in the model, and they agree with the Key Finding of each dimension. Each dimension has one metric that behaves like the confidence (deletion AUC, mask-out drop, margin after removal, necessity \(F^{-}\)) and one that describes the explainer (insertion AUC, sufficiency drop, margin of the explanation alone, sufficiency \(F^{+}\)).

In the combined detector the first group is redundant with the confidence and loses its weight: on AG News each of these metrics is significant in at most four of the eight graph configurations, and its sign often changes from one topology to the next. The second group keeps its weight: the margin of the explanation alone is significant in seven of the eight configurations and the insertion AUC in six, with the same sign throughout. On SST-2, where the explanation adds nothing, no feature is significant in more than four configurations and TokenSHAP has only three small significant coefficients (Table~S4).

In short, whatever the combined detector gains comes from the weaker error signals of each dimension, the ones that describe the explainer, because the stronger ones are already in the model as the confidence; and, as shown above, that gain never reaches the confidence of the base model.

\paragraph{\textbf{Key Finding}}
Explanation signatures from any explainer and representation do not detect classification errors beyond the confidence of the language model, and no graph configuration detects them better than the token explainer of the language model itself. What the explanation does add is limited to the distilled surrogate: it is over-confident in the teacher's errors, which its confidence therefore hides, and the perturbation-based metrics recover part of that hidden discrimination, without ever reaching the confidence of the base model.

\subsection{What the Surrogate Inherits from Its Teacher} \label{sec:calibration} \label{sec:ablation}

The surrogate receives two things from its teacher, the labels it is trained on and the node features of its graphs, both outputs of the same fine-tuned checkpoint (Section~\ref{gnn}). Through them it could inherit the teacher's decisions and errors, become confident in them, and inherit its attributions. We test each possibility with its own experiment: calibration against the teacher's and the gold labels, a control surrogate trained on the gold labels, and an ablation of the node features.

\subsubsection{Does the surrogate inherit the teacher's confidence?} \label{sec:inherit_conf}
Table~\ref{tab:calibration} reports, for each surrogate and for the teacher, the expected calibration error (ECE, Section~\ref{logistic_regression}) and the number of disagreements or errors, with respect to the labels the surrogate was trained on and with respect to the gold labels. Beyond how many errors a model makes, we want to know how sure it is when it makes them, so the table also counts the errors made with confidence above 0.9.

\begin{table}[H]
\centering
\scriptsize
\setlength{\tabcolsep}{4pt}
\caption{Calibration of each surrogate (one per graph type, shared by the two explainers) and of the teacher, with respect to the teacher's labels and to the gold labels: expected calibration error (ECE), disagreements or errors (fraction of the test set), and how many of them are made with confidence above 0.9 (fraction of the count); --: not applicable to the teacher. Error detection among the high-confidence predictions in Table~S13.}
\label{tab:calibration}
\begin{tabular}{@{}llcccccc@{}}
\toprule
& & \multicolumn{3}{c}{\textbf{With respect to the teacher's labels}} & \multicolumn{3}{c}{\textbf{With respect to the gold labels}} \\
\cmidrule(lr){3-5} \cmidrule(lr){6-8}
\textbf{Dataset} & \textbf{Model} & ECE & Disagreements & of which conf.\ \(>0.9\) & ECE & Errors & of which conf.\ \(>0.9\) \\
\midrule
\multirow{4}{*}{SubgraphX} & Constituency & -- & 574 & 0.50 & 288 & 0.777 / 0.769 & \(-0.008\) [\(-0.028, +0.011\)] & \(+0.003\) [\(-0.017, +0.023\)] \\
 & Syntactic & -- & 606 & 0.44 & 268 & 0.759 / 0.766 & \(+0.007\) [\(-0.015, +0.027\)] & \(+0.007\) [\(-0.025, +0.041\)] \\
 & Window & -- & 562 & 0.49 & 278 & 0.771 / 0.787 & \(+0.016\) [\(-0.005, +0.038\)] & \(+0.014\) [\(-0.003, +0.033\)] \\
 & Skip-gram & -- & 564 & 0.50 & 282 & 0.784 / 0.788 & \(+0.005\) [\(-0.019, +0.029\)] & \(+0.008\) [\(-0.005, +0.022\)] \\
\midrule
\multirow{4}{*}{GraphSVX} & Constituency & -- & 574 & 0.50 & 288 & 0.777 / 0.775 & \(-0.002\) [\(-0.026, +0.020\)] & \(+0.000\) [\(-0.015, +0.016\)] \\
 & Syntactic & -- & 606 & 0.44 & 268 & 0.759 / 0.775 & \(+0.016\) [\(-0.008, +0.038\)] & \(-0.008\) [\(-0.042, +0.026\)] \\
 & Window & -- & 562 & 0.49 & 278 & 0.770 / 0.782 & \(+0.012\) [\(-0.012, +0.035\)] & \(+0.025\) [\(+0.005, +0.050\)] \\
 & Skip-gram & -- & 564 & 0.50 & 282 & 0.779 / 0.800 & \(+0.022\) [\(+0.000, +0.043\)] & \(+0.026\) [\(+0.010, +0.049\)] \\
\midrule
TokenSHAP & Tokens & -- & 471 & 0.31 & 144 & 0.809 / 0.795 & \(-0.013\) [\(-0.029, +0.004\)] & \(+0.000\) [\(-0.017, +0.021\)] \\

\end{tabular}
\end{table}

The surrogates are well calibrated with respect to the labels they were trained on, the teacher's: they disagree with it rarely and are sure of only 6--12\% of their disagreements (Table~\ref{tab:calibration}), which is why disagreement with the teacher is detected from the confidence alone (Table~S5) and the explanation adds nothing to it. With respect to the gold labels they are four times worse calibrated than BERT on AG News, and seven to fourteen times worse than with respect to the teacher's labels (two to four times on SST-2): their confidence tracks their agreement with the teacher, not the truth, and the teacher's errors, which BERT makes hesitantly, the surrogate makes with confidence. On AG News the surrogates make as many errors as BERT but are sure of twice as many of them; on SST-2 the gap to BERT is small because BERT itself is sure of almost half of its errors. This is why the surrogate's confidence is a weaker error signal than BERT's on AG News and not on SST-2, according to the asymmetry of Figure~\ref{fig:error_detection}. Within those high-confidence predictions, where confidence alone can no longer tell an error from a correct prediction, the explanation still raises the AUROC of the surrogate's detector in every AG News configuration (significantly in five of the eight) and in none of the SST-2 configurations beyond the bootstrap interval (Table~S13): this is where the gain of Figure~\ref{fig:error_detection} comes from. The reference for these gains is the same errors, not the language model's own regime, which holds different ones: adding the surrogate's signals to BERT's detector among BERT's high-confidence predictions gives no gain (Table~S7), and on the whole test set the surrogate's combined detector remains below BERT's confidence (Figure~\ref{fig:error_detection}).

\subsubsection{Does it inherit the errors, and the attributions?} \label{sec:inherit_err}
To separate what comes from the teacher's labels from what comes from the instances, every surrogate was retrained on the gold labels with the same graphs, the same node features from the fine-tuned teacher, and the same architecture and training protocol, and once more on the teacher labels with a different seed (Section~\ref{logistic_regression}); Table~\ref{tab:control} compares both against BERT's errors.

\begin{table}[H]
\centering
\scriptsize
\setlength{\tabcolsep}{4pt}
\caption{Gold-trained control. For each topology, the distilled surrogate (trained on the teacher's labels) and the gold-trained surrogate (same graphs, node features, architecture and protocol, trained on the gold labels) are compared against BERT's gold-label errors on the test split: accuracy with respect to the gold labels; fraction of BERT's errors on which the surrogate predicts the same wrong class as BERT; fraction of BERT's errors on which it does so with confidence above 0.9. \(^{*}\): the 95\% paired-bootstrap interval of the difference between the two surrogates excludes zero (intervals in Table~S14). A seed replicate of each surrogate does not narrow the gap (Table~S9).}
\label{tab:control}
\resizebox{\textwidth}{!}{%
\begin{tabular}{@{}llccc@{}}
\toprule
\textbf{Dataset} & \textbf{Graph} & \textbf{Accuracy (gold)} & \textbf{Same wrong class as BERT on its errors} & \textbf{\ldots and confident (\(>0.9\))} \\
& & Distilled / Gold-trained & Distilled / Gold-trained & Distilled / Gold-trained \\
\midrule
\multirow{4}{*}{AG News} & Constituency & 0.924 / 0.926 & 0.84 / 0.80 & \(-0.036\) [\(-0.073, -0.002\)] & 0.55 / 0.27 & \(-0.276\) [\(-0.322, -0.230\)] \\
 & Syntactic & 0.920 / 0.924 & 0.82 / 0.74 & \(-0.084\) [\(-0.120, -0.046\)] & 0.50 / 0.22 & \(-0.278\) [\(-0.322, -0.232\)] \\
 & Window & 0.926 / 0.922 & 0.83 / 0.80 & \(-0.031\) [\(-0.070, +0.004\)] & 0.52 / 0.24 & \(-0.282\) [\(-0.327, -0.242\)] \\
 & Skip-gram & 0.926 / 0.931 & 0.84 / 0.75 & \(-0.082\) [\(-0.121, -0.044\)] & 0.54 / 0.23 & \(-0.304\) [\(-0.350, -0.261\)] \\
\midrule
\multirow{4}{*}{SST-2} & Constituency & 0.890 / 0.899 & 0.71 / 0.66 & \(-0.053\) [\(-0.145, +0.031\)] & 0.41 / 0.30 & \(-0.109\) [\(-0.198, -0.025\)] \\
 & Syntactic & 0.875 / 0.859 & 0.77 / 0.74 & \(-0.025\) [\(-0.127, +0.080\)] & 0.41 / 0.19 & \(-0.218\) [\(-0.333, -0.118\)] \\
 & Window & 0.888 / 0.893 & 0.75 / 0.70 & \(-0.057\) [\(-0.149, +0.027\)] & 0.41 / 0.25 & \(-0.165\) [\(-0.271, -0.078\)] \\
 & Skip-gram & 0.878 / 0.893 & 0.71 / 0.67 & \(-0.040\) [\(-0.143, +0.067\)] & 0.42 / 0.36 & \(-0.067\) [\(-0.157, +0.028\)] \\

\end{tabular}%
}
\end{table}

A gold-trained surrogate makes almost all of the same errors as BERT, with the same wrong class: on 84--92\% of BERT's errors, against 85--92\% for the distilled surrogate, a difference of at most 3 points that is not significant in any topology (Table~S14), and the two surrogates make their errors on the same instances (Jaccard 0.76--0.89, against 0.03--0.05 if their errors were independent). The errors are therefore a property of the instances and of the node features, not of the teacher's labels. What distillation adds is confidence in them: the distilled surrogate gives BERT's wrong class with confidence above 0.9 on 51--63\% of BERT's errors, the gold-trained one on 25--38\% (significant in all eight topologies); and the distilled surrogate is neither a better classifier (accuracy within 0.4 points) nor more faithful to BERT's predictions. The attributions are not inherited either: the words highlighted by the graph explainers overlap with those highlighted by TokenSHAP at or just above the level expected by chance (Jaccard 0.08--0.18, against 0.11 for two independent selections of the same size; rank correlation 0.01--0.19), and the overlap does not predict the error of either model (Table~S8).

\subsubsection{Through which channel do the errors arrive?} \label{sec:inherit_channel}
Every result so far uses node features from the fine-tuned teacher. Repeating the whole pipeline with node features from the pre-trained encoder (\texttt{bert-base-uncased}, everything else unchanged, gold-trained controls included) ablates the one ingredient that distinguishes a task-adapted from a generic representation (Table~\ref{tab:ablation}; full tables in S12--S14).

Holding the teacher's labels fixed and replacing the fine-tuned features by pre-trained ones lowers the fraction of BERT's errors on which the surrogate gives BERT's wrong class from 85--92\% to 71--82\%; holding the fine-tuned features fixed and replacing the teacher's labels by the gold labels changes it by at most 3 points, and the twins still make their errors on the same instances. Only with pre-trained features do the labels carry part of the errors (the distilled surrogate gives BERT's wrong class 3--9 points more often than its gold-trained twin, and the error sets of the twins overlap less, Jaccard 0.53--0.66).

The errors travel through the node features, and the labels matter only when the features do not already encode the teacher's decision.

\begin{table}[H]
\centering
\scriptsize
\setlength{\tabcolsep}{3pt}
\caption{Ablation of the node features (pre-trained encoder against the fine-tuned teacher). Agreement, accuracy and ECE: medians over the four topologies; other rows: ranges over the eight configurations. Error detection: AUROC of the gold-error detector on the surrogate's confidence and on confidence plus explanation, and configurations whose nested gain has a 95\% interval above zero. Last block: fraction of BERT's errors on which the surrogate gives BERT's wrong class, and with confidence above 0.9 (distilled / gold-trained). Full tables in S12--S14.}
\label{tab:ablation}
\begin{tabular}{@{}lcccc@{}}
\toprule
& \multicolumn{2}{c}{\textbf{AG News}} & \multicolumn{2}{c}{\textbf{SST-2}} \\
\cmidrule(lr){2-3} \cmidrule(lr){4-5}
& pre-trained & fine-tuned & pre-trained & fine-tuned \\
\midrule
\emph{Surrogate quality} & & & & \\
\quad Agreement with the teacher & 0.96 & 0.98 & 0.92 & 0.98 \\
\quad Accuracy (gold labels) & 0.93 & 0.94 & 0.88 & 0.91 \\
\addlinespace
\emph{Calibration (gold labels)} & & & & \\
\quad ECE & 0.042 & 0.042 & 0.047 & 0.052 \\
\quad Errors with confidence \(>0.9\) & 44--50\% & 56--61\% & 33--41\% & 49--57\% \\
\addlinespace
\emph{Error detection (AUROC)} & & & & \\
\quad Surrogate confidence & 0.84--0.85 & 0.82--0.83 & 0.80--0.82 & 0.83--0.85 \\
\quad Confidence + explanation & 0.83--0.85 & 0.83--0.86 & 0.77--0.81 & 0.80--0.84 \\
\quad Configurations with a significant gain (of 8) & 0 & 7 & 0 & 0 \\
\addlinespace
\emph{On BERT's errors (distilled / gold-trained)} & & & & \\
\quad Same wrong class as BERT & 79--82 / 70--76 & 85--88 / 84--85 & 71--77 / 66--74 & 89--92 / 89--92 \\
\quad \ldots and confident (\(>0.9\)) & 50--55 / 22--27 & 55--63 / 25--31 & 41--42 / 19--36 & 51--58 / 26--38 \\

\end{tabular}
\end{table}

The ablation also shows that the conclusions do not depend on the features. Fine-tuned features make the surrogates more faithful and more accurate (Table~\ref{tab:ablation}), and change the route, not the amount, of the error information that reaches the detector: with pre-trained features the surrogate is less over-confident in its errors (44--50\% of its errors above confidence 0.9 against 56--61\%) and the explanation adds nothing to its confidence; with fine-tuned features the confidence loses discrimination and the explanation recovers part of it (Section~\ref{sec:logistic}).

The combined detector ends at the same level in both settings, no surrogate signal improves on BERT's own confidence in either, and the gold-trained control separates the errors from the confidence with either representation.

\paragraph{\textbf{Key Finding}}
Distillation transfers the teacher's decisions and makes the surrogate confident in them, but transfers neither its errors nor its reasoning: only the distilled surrogate is confident in the teacher's errors (Section~\ref{sec:inherit_conf}); a gold-trained twin makes its errors on the same instances, and the attributions of the two models overlap at chance level (Section~\ref{sec:inherit_err}); and the errors reach the surrogate through the node features, since replacing them moves the overlap with BERT's errors and replacing the labels does not (Section~\ref{sec:inherit_channel}).

\paragraph{\textbf{The results read together}}
Four dimensions, two datasets, three explainers and four graph topologies return one answer: the prediction rests on far more evidence than the explanation contains. Removing the 20\% of the input that the explainer selects leaves the confidence almost intact (Dimensions~1, 2 and~4) and changes the decision only where the margin was already narrow (Dimension~3). Two consequences follow.

First, every metric that measures what happens after the explanation is removed is the model's confidence in another form, and it is the informative metric of its dimension (Table~S11); the metrics that describe the explanation alone are weaker error signals and differ between explainers, not between representations. An explanation-based detector can therefore repeat what the confidence already says, but not exceed it: no explanation signature from any explainer or representation detects errors beyond the confidence of the language model (Figure~\ref{fig:error_detection}).

Second, the one gain observed, that of the explanation over a distilled surrogate's own confidence, is a property of distillation, not of graphs, and its meaning is precise: it brings the surrogate closer to its base model. The surrogate inherits the teacher's decisions and is confident in them, including the teacher's errors; the perturbation-based metrics recover part of what that confidence hides, and each part recovered moves the surrogate's error detection towards the level of the language model without ever reaching it (Figure~\ref{fig:error_detection}). The errors themselves belong to the instances and their features, since a gold-trained twin makes them too, and the surrogate does not inherit the teacher's reasoning: its attributions overlap with the teacher's at chance level. The one difference between the two families, the language model's stronger dependence on its top-ranked tokens, is not an advantage of the graphs: at no point does a graph surrogate, with or without its explanation, reach the language model, neither in accuracy nor in error detection.
